\subsection{Dictionary Construction Patterns}

A dictionary can be constructed in several ways, but all construction patterns produce the same kind of object: a \texttt{dict}. The differences are mostly about where the key-value pairs come from.

The common construction patterns are:

\begin{itemize}
    \item writing key-value pairs directly with a dictionary literal;
    \item calling the \texttt{dict()} constructor;
    \item generating pairs with a dictionary comprehension;
    \item converting a sequence of two-item pairs into a dictionary.
\end{itemize}

All of these patterns must still obey the same dictionary rules. Keys must be hashable. Equal keys collapse into one mapping entry. Values may be arbitrary Python objects.

\subsubsection{Literal Construction with \texttt{\{key: value\}} Pairs}

The most direct way to build a dictionary is the literal form. Each entry is written as a key, a colon, and a value.

\begin{verbatim}
quote_by_name = {
    "Homer": "D'oh!",
    "Jerry": "No soup for you!",
    "Arthur": "It is only a model.",
}

print(f"quote_by_name: {quote_by_name!r}")
print(f"type(quote_by_name): {type(quote_by_name)}")
print(f"len(quote_by_name): {len(quote_by_name)}")
\end{verbatim}

Possible output:

\begin{verbatim}
quote_by_name: {'Homer': "D'oh!", 'Jerry': 'No soup for you!', 'Arthur': 'It is only a model.'}
type(quote_by_name): <class 'dict'>
len(quote_by_name): 3
\end{verbatim}

The dictionary has three keys and three associated values.

A few important dictionary methods can be inspected with \texttt{dir()}.

\begin{verbatim}
quote_by_name = {
    "Homer": "D'oh!",
    "Jerry": "No soup for you!",
    "Arthur": "It is only a model.",
}

print(f"type(quote_by_name): {type(quote_by_name)}")
print(f"__getitem__ in dir: {'__getitem__' in dir(quote_by_name)}")
print(f"__setitem__ in dir: {'__setitem__' in dir(quote_by_name)}")
print(f"keys in dir:        {'keys' in dir(quote_by_name)}")
print(f"values in dir:      {'values' in dir(quote_by_name)}")
print(f"items in dir:       {'items' in dir(quote_by_name)}")
print(f"update in dir:      {'update' in dir(quote_by_name)}")
\end{verbatim}

Possible output:

\begin{verbatim}
type(quote_by_name): <class 'dict'>
__getitem__ in dir: True
__setitem__ in dir: True
keys in dir:        True
values in dir:      True
items in dir:       True
update in dir:      True
\end{verbatim}

The central methods here are:

\begin{itemize}
    \item \texttt{\_\_getitem\_\_}: value retrieval by key;
    \item \texttt{\_\_setitem\_\_}: key-value insertion or replacement;
    \item \texttt{keys}: view over the key domain;
    \item \texttt{values}: view over the value domain;
    \item \texttt{items}: view over key-value pairs;
    \item \texttt{update}: later used to merge or overwrite mappings.
\end{itemize}

Literal construction can use different hashable key types.

\begin{verbatim}
mixed_keys = {
    "answer": 42,
    42: "the number itself",
    ("x", "y"): "coordinate names",
}

print(f"mixed_keys: {mixed_keys!r}")
print(f"mixed_keys['answer']: {mixed_keys['answer']!r}")
print(f"mixed_keys[42]: {mixed_keys[42]!r}")
print(f"mixed_keys[('x', 'y')]: {mixed_keys[('x', 'y')]!r}")
\end{verbatim}

Possible output:

\begin{verbatim}
mixed_keys: {'answer': 42, 42: 'the number itself', ('x', 'y'): 'coordinate names'}
mixed_keys['answer']: 42
mixed_keys[42]: 'the number itself'
mixed_keys[('x', 'y')]: 'coordinate names'
\end{verbatim}

The keys are not required to be strings. They are required to be hashable.

If the same key appears more than once in a literal, the later value wins.

\begin{verbatim}
data = {
    "name": "Homer",
    "job": "safety inspector",
    "job": "donut inspector",
}

print(f"data: {data!r}")
print(f"len(data): {len(data)}")
print(f"data['job']: {data['job']!r}")
\end{verbatim}

Possible output:

\begin{verbatim}
data: {'name': 'Homer', 'job': 'donut inspector'}
len(data): 2
data['job']: 'donut inspector'
\end{verbatim}

The dictionary does not store two separate \texttt{"job"} entries. The second \texttt{"job"} assignment replaces the earlier value for the same key.

\subsubsection{Constructor-Based Construction with \texttt{dict()}}

The \texttt{dict()} constructor creates dictionaries from several kinds of input.

With no arguments, it creates an empty dictionary.

\begin{verbatim}
empty_data = dict()

print(f"empty_data: {empty_data!r}")
print(f"type(empty_data): {type(empty_data)}")
print(f"len(empty_data): {len(empty_data)}")
\end{verbatim}

Possible output:

\begin{verbatim}
empty_data: {}
type(empty_data): <class 'dict'>
len(empty_data): 0
\end{verbatim}

This is equivalent in result to \texttt{\{\}}, although \texttt{\{\}} is more common for an empty dictionary.

The constructor can also use keyword arguments.

\begin{verbatim}
person = dict(name="Homer", job="safety inspector", age=42)

print(f"person: {person!r}")
print(f"type(person): {type(person)}")
print(f"person['name']: {person['name']!r}")
\end{verbatim}

Possible output:

\begin{verbatim}
person: {'name': 'Homer', 'job': 'safety inspector', 'age': 42}
type(person): <class 'dict'>
person['name']: 'Homer'
\end{verbatim}

In this form, the keyword names become string keys. The expression above creates keys \texttt{"name"}, \texttt{"job"}, and \texttt{"age"}.

This form is limited to keys that are valid Python identifier names. It cannot directly create keys such as \texttt{"first name"}, \texttt{42}, or \texttt{("x", "y")}.

\begin{verbatim}
literal_data = {
    "first name": "Homer",
    42: "answer",
    ("x", "y"): "coordinate names",
}

print(f"literal_data: {literal_data!r}")
\end{verbatim}

Possible output:

\begin{verbatim}
literal_data: {'first name': 'Homer', 42: 'answer', ('x', 'y'): 'coordinate names'}
\end{verbatim}

For those keys, the literal form is clearer and more general.

The constructor can also copy an existing dictionary.

\begin{verbatim}
original = {
    "name": "Marge",
    "quote": "Mmm-hmm.",
}

copy_data = dict(original)

print(f"original:  {original!r}")
print(f"copy_data: {copy_data!r}")
print(f"same object: {original is copy_data}")
print(f"equal value: {original == copy_data}")
\end{verbatim}

Possible output:

\begin{verbatim}
original:  {'name': 'Marge', 'quote': 'Mmm-hmm.'}
copy_data: {'name': 'Marge', 'quote': 'Mmm-hmm.'}
same object: False
equal value: True
\end{verbatim}

The new dictionary has the same key-value pairs, but it is a different dictionary object.

This is a shallow copy. If a value is a mutable object, the value object itself is still shared.

\begin{verbatim}
original = {
    "name": "Homer",
    "quotes": ["D'oh!"],
}

copy_data = dict(original)

copy_data["quotes"].append("Mmm... donuts")

print(f"original:  {original!r}")
print(f"copy_data: {copy_data!r}")
print(f"same list object: {original['quotes'] is copy_data['quotes']}")
\end{verbatim}

Possible output:

\begin{verbatim}
original:  {'name': 'Homer', 'quotes': ["D'oh!", 'Mmm... donuts']}
copy_data: {'name': 'Homer', 'quotes': ["D'oh!", 'Mmm... donuts']}
same list object: True
\end{verbatim}

The dictionary object was copied, but the nested list was not duplicated.

\subsubsection{Dictionary Comprehensions as Key-Value Generation Pipelines}

A dictionary comprehension builds a dictionary by generating key-value pairs from an iterable source.

The general shape is:

\begin{verbatim}
{key_expression: value_expression for item in source}
\end{verbatim}

A simple example maps names to their lengths.

\begin{verbatim}
names = ["Homer", "Marge", "Lisa", "Bart"]

name_lengths = {name: len(name) for name in names}

print(f"names: {names!r}")
print(f"name_lengths: {name_lengths!r}")
print(f"type(name_lengths): {type(name_lengths)}")
\end{verbatim}

Possible output:

\begin{verbatim}
names: ['Homer', 'Marge', 'Lisa', 'Bart']
name_lengths: {'Homer': 5, 'Marge': 5, 'Lisa': 4, 'Bart': 4}
type(name_lengths): <class 'dict'>
\end{verbatim}

The expression before the colon generates the key. The expression after the colon generates the value.

The same construction can be written more mechanically as a loop.

\begin{verbatim}
names = ["Homer", "Marge", "Lisa", "Bart"]

name_lengths = {}

for name in names:
    name_lengths[name] = len(name)

print(f"name_lengths: {name_lengths!r}")
\end{verbatim}

Possible output:

\begin{verbatim}
name_lengths: {'Homer': 5, 'Marge': 5, 'Lisa': 4, 'Bart': 4}
\end{verbatim}

The comprehension is the compact form of this construction pattern.

A dictionary comprehension can include a filter.

\begin{verbatim}
names = ["Homer", "Marge", "Lisa", "Bart", "Maggie"]

long_names = {name: len(name) for name in names if len(name) > 4}

print(f"names: {names!r}")
print(f"long_names: {long_names!r}")
\end{verbatim}

Possible output:

\begin{verbatim}
names: ['Homer', 'Marge', 'Lisa', 'Bart', 'Maggie']
long_names: {'Homer': 5, 'Marge': 5, 'Maggie': 6}
\end{verbatim}

Only names with length greater than \texttt{4} become dictionary entries.

The generated key must still be hashable.

\begin{verbatim}
words = ["spam", "eggs"]

try:
    bad_map = {[word]: len(word) for word in words}
except TypeError as err:
    print(f"comprehension failed: {err}")
\end{verbatim}

Possible output:

\begin{verbatim}
comprehension failed: unhashable type: 'list'
\end{verbatim}

The generated key \texttt{[word]} is a list. Lists are unhashable, so the dictionary cannot use them as keys.

If the comprehension generates the same key more than once, the later value replaces the earlier one.

\begin{verbatim}
words = ["SPAM", "spam", "Spam"]

normalized = {word.lower(): word for word in words}

print(f"words: {words!r}")
print(f"normalized: {normalized!r}")
print(f"len(normalized): {len(normalized)}")
\end{verbatim}

Possible output:

\begin{verbatim}
words: ['SPAM', 'spam', 'Spam']
normalized: {'spam': 'Spam'}
len(normalized): 1
\end{verbatim}

All three source words generate the same key, \texttt{"spam"}. The last generated value survives.

A dictionary comprehension is useful when the program naturally computes a mapping.

\begin{verbatim}
quotes = [
    "D'oh!",
    "No soup for you!",
    "It is only a model.",
]

quote_lengths = {quote: len(quote) for quote in quotes}

for quote, size in quote_lengths.items():
    print(f"{quote!r:<28} -> {size}")
\end{verbatim}

Possible output:

\begin{verbatim}
"D'oh!"                      -> 5
'No soup for you!'           -> 16
'It is only a model.'        -> 19
\end{verbatim}

The result maps each quote to a computed property of that quote.

\subsubsection{Building Dictionaries from Pair Sequences}

The \texttt{dict()} constructor can also build a dictionary from a sequence of pairs. Each pair must contain exactly two elements: the key and the value.

\begin{verbatim}
pairs = [
    ("Homer", "D'oh!"),
    ("Jerry", "No soup for you!"),
    ("Arthur", "It is only a model."),
]

quote_by_name = dict(pairs)

print(f"pairs: {pairs!r}")
print(f"type(pairs): {type(pairs)}")
print()
print(f"quote_by_name: {quote_by_name!r}")
print(f"type(quote_by_name): {type(quote_by_name)}")
\end{verbatim}

Possible output:

\begin{verbatim}
pairs: [('Homer', "D'oh!"), ('Jerry', 'No soup for you!'), ('Arthur', 'It is only a model.')]
type(pairs): <class 'list'>

quote_by_name: {'Homer': "D'oh!", 'Jerry': 'No soup for you!', 'Arthur': 'It is only a model.'}
type(quote_by_name): <class 'dict'>
\end{verbatim}

The outer object is a list. Each inner object is a two-item tuple. The constructor reads each pair as:

\begin{verbatim}
(first item)  -> key
(second item) -> value
\end{verbatim}

The inner pair type can be another iterable with two elements.

\begin{verbatim}
pairs = [
    ["Homer", "D'oh!"],
    ["Jerry", "No soup for you!"],
]

quote_by_name = dict(pairs)

print(f"quote_by_name: {quote_by_name!r}")
\end{verbatim}

Possible output:

\begin{verbatim}
quote_by_name: {'Homer': "D'oh!", 'Jerry': 'No soup for you!'}
\end{verbatim}

The inner lists are accepted as pair containers. The keys themselves are strings, so they are hashable.

However, each pair must have exactly two elements.

\begin{verbatim}
bad_pairs = [
    ("Homer", "D'oh!"),
    ("Jerry", "No soup for you!", "extra"),
]

try:
    data = dict(bad_pairs)
except ValueError as err:
    print(f"construction failed: {err}")
\end{verbatim}

Possible output:

\begin{verbatim}
construction failed: dictionary update sequence element #1 has length 3; 2 is required
\end{verbatim}

The second pair has three elements, so Python cannot interpret it as one key and one value.

Duplicate keys collapse during construction.

\begin{verbatim}
pairs = [
    ("Homer", "D'oh!"),
    ("Homer", "Mmm... donuts"),
    ("Lisa", "If anyone wants me, I'll be in my room."),
]

data = dict(pairs)

print(f"data: {data!r}")
print(f"len(data): {len(data)}")
print(f"data['Homer']: {data['Homer']!r}")
\end{verbatim}

Possible output:

\begin{verbatim}
data: {'Homer': 'Mmm... donuts', 'Lisa': "If anyone wants me, I'll be in my room."}
len(data): 2
data['Homer']: 'Mmm... donuts'
\end{verbatim}

The later \texttt{"Homer"} pair overwrote the earlier \texttt{"Homer"} pair.

A common pair-producing function is \texttt{zip()}. It combines parallel iterables into pairs.

\begin{verbatim}
names = ["Homer", "Jerry", "Arthur"]
quotes = ["D'oh!", "No soup for you!", "It is only a model."]

pairs = zip(names, quotes)
quote_by_name = dict(pairs)

print(f"names: {names!r}")
print(f"quotes: {quotes!r}")
print(f"quote_by_name: {quote_by_name!r}")
\end{verbatim}

Possible output:

\begin{verbatim}
names: ['Homer', 'Jerry', 'Arthur']
quotes: ["D'oh!", 'No soup for you!', 'It is only a model.']
quote_by_name: {'Homer': "D'oh!", 'Jerry': 'No soup for you!', 'Arthur': 'It is only a model.'}
\end{verbatim}

The conceptual transformation is:

\begin{verbatim}
"Homer"  + "D'oh!"                -> ("Homer", "D'oh!")
"Jerry"  + "No soup for you!"     -> ("Jerry", "No soup for you!")
"Arthur" + "It is only a model."  -> ("Arthur", "It is only a model.")
\end{verbatim}

Then \texttt{dict()} turns those pairs into key-value entries.

The construction patterns can be summarized as follows:

\begin{center}
\begin{tabular}{ll}
\textbf{Construction Form} & \textbf{Typical Use} \\
\hline
\texttt{\{"a": 1, "b": 2\}} & Directly known key-value pairs \\
\texttt{dict()} & Empty dictionary or constructor-based creation \\
\texttt{dict(name="Homer")} & Simple string keys that are valid identifiers \\
\texttt{\{k: v for ...\}} & Generated mappings \\
\texttt{dict([("a", 1), ("b", 2)])} & Pair sequence conversion \\
\texttt{dict(zip(keys, values))} & Parallel key and value sources \\
\end{tabular}
\end{center}

All forms create the same object type: \texttt{dict}. The difference is the source and shape of the input data.